\chapter{Instrukcja eksportu Mermaid do LaTeX}
\label{chap:mermaid-latex}

\section{Cel dodatku}
W dodatku przedstawiono powtarzalną procedurę przekształcenia diagramów Mermaid (C4) do postaci grafik osadzanych w dokumencie LaTeX.

\section{Organizacja plików}
Zalecana struktura:
\begin{itemize}
    \item źródła Mermaid: \texttt{images/diagrams-src/*.mmd},
    \item wygenerowane grafiki: \texttt{images/diagrams/*.pdf} lub \texttt{images/diagrams/*.svg}.
\end{itemize}

Rekomendowane nazwy plików:
\begin{itemize}
    \item \texttt{c1-context},
    \item \texttt{c2-container},
    \item \texttt{c3-component-smf},
    \item \texttt{c4-deployment}.
\end{itemize}

\section{Eksport z Mermaid CLI}
Do eksportu można użyć narzędzia \texttt{mmdc}:
\begin{lstlisting}[basicstyle=\footnotesize\ttfamily,caption=Eksport diagramu do SVG,label=lst:mmdc-svg]
mmdc -i diagram.mmd -o images/diagram.svg -t neutral -b transparent
\end{lstlisting}

\begin{lstlisting}[basicstyle=\footnotesize\ttfamily,caption=Eksport diagramu do PDF,label=lst:mmdc-pdf]
mmdc -i diagram.mmd -o images/diagram.pdf -t neutral -b white
\end{lstlisting}

\section{Osadzanie diagramu w rozdziale LaTeX}
Przykład osadzenia:
\begin{lstlisting}[basicstyle=\footnotesize\ttfamily,caption=Wstawienie diagramu C2 do LaTeX,label=lst:latex-include-c2]
\begin{figure}[htbp]
  \centering
  \includegraphics[width=0.95\linewidth]{images/diagrams/c2-container.pdf}
  \caption{C2 Container Diagram dla Shadow Deploy SMF}
  \label{fig:c2-shadow-smf}
\end{figure}
\end{lstlisting}

\section{Jak wyodrębnić kod Mermaid}
Kod diagramu należy skopiować z bloku oznaczonego:
\begin{itemize}
    \item początkiem \texttt{```mermaid},
    \item końcem \texttt{```}.
\end{itemize}

Każdy diagram powinien zostać zapisany do osobnego pliku \texttt{.mmd}, a następnie wyeksportowany do PDF lub SVG.

\section{Dobre praktyki}
\begin{itemize}
    \item Do wersji drukowanej preferować PDF.
    \item Stosować spójne nazwy plików i etykiet \texttt{\textbackslash label\{\}}.
    \item Po każdej zmianie architektury regenerować grafiki z tych samych plików źródłowych.
\end{itemize}